\section{Working alongside the analyst}

The unit of work between an agent and the analyst is a \emph{suggestion}: a proposed change to a set of cells, such as the capacities of a new pipeline or the reference price of a node.
Figure \ref{fig:review} shows how a suggestion could look on the analyst's screen and what happens to the analyst's answer.

\begin{figure}[htbp]
\centering
% Figure 5.5: one suggestion on the analyst's screen, the analyst's reply, and where the reply goes.
\begin{tikzpicture}[ggm, x=1cm, y=1cm]
  \tikzset{ph/.style={line width=1.8pt, draw=ggmMuted!40, line cap=round},
           lblcap/.style={font=\sffamily\tiny\bfseries, text=ggmMuted, anchor=west, inner sep=0.5mm},
           lblbody/.style={font=\sffamily\scriptsize, text=ggmInk, anchor=west, inner sep=0.5mm}}
  % ---------- the monitor ----------
  \filldraw[fill=ggmModel, draw=ggmModel, rounded corners=2mm] (0,3.7) rectangle (7.75,10.9);
  \fill[white] (0.14,3.84) rectangle (7.61,10.76);
  \fill[ggmModel] (3.5,3.28) -- (4.25,3.28) -- (4.1,3.7) -- (3.65,3.7) -- cycle;
  \fill[ggmModel!70, rounded corners=0.6mm] (2.9,3.15) rectangle (4.85,3.3);
  \fill[ggmMuted!20] (0.14,10.46) rectangle (7.61,10.76);
  \foreach \x/\c in {0.34/ggmAnalyst, 0.52/ggmSource, 0.7/ggmRecord} \fill[\c] (\x,10.61) circle (0.06);
  \node[font=\sffamily\scriptsize, text=ggmModel] at (3.875,10.61) {GGM workbench \textperiodcentered{} data step};
  % --- spreadsheet pane: columns at 0.3 | 1.32 | 2.47 | 3.55, nine rows of 0.45 ---
  \fill[ggmSheet] (0.3,9.95) rectangle (3.55,10.3);
  \node[font=\sffamily\scriptsize\bfseries, text=white, anchor=west] at (0.34,10.125) {Sheet: Nodes};
  \fill[ggmSheet!12] (0.3,9.6) rectangle (3.55,9.95);
  \node[lblbody] at (0.38,9.775) {Node};
  \node[lblbody] at (1.38,9.775) {Price};
  \node[lblbody] at (2.53,9.775) {Conf.};
  \fill[ggmAgent!14] (0.3,7.8) rectangle (3.55,8.25);
  \foreach \k/\c/\len in {0/ggmRecord/0.55, 1/ggmRecord/0.65, 2/ggmSource/0.75, 4/ggmRecord/0.65,
                          5/ggmAnalyst/0.75, 6/ggmRecord/0.55, 7/red!65!black/0.65, 8/ggmSource/0.75} {
    \draw[ph] (0.45,9.375-\k*0.45) -- ++(\len,0);
    \draw[ph] (1.47,9.375-\k*0.45) -- (2.15,9.375-\k*0.45);
    \fill[\c] (3.0,9.375-\k*0.45) circle (0.08);
  }
  \draw[ph] (0.45,8.025) -- ++(0.65,0);
  \draw[ggmSheet!35] (0.3,5.55) rectangle (3.55,9.6);
  \foreach \x in {1.32,2.47} \draw[ggmSheet!35] (\x,5.55) -- (\x,9.95);
  \foreach \r in {1,...,8} \draw[ggmSheet!22] (0.3,9.6-\r*0.45) -- (3.55,9.6-\r*0.45);
  % the proposed cell
  \filldraw[fill=white, draw=ggmAgent, line width=1pt, rounded corners=0.6mm] (1.38,7.86) rectangle (2.41,8.19);
  \draw[line width=1.8pt, draw=ggmAgent!70, line cap=round] (1.5,8.025) -- (2.1,8.025);
  \node[font=\sffamily\footnotesize\bfseries, text=ggmAgent] at (3.0,8.025) {?};
  % confidence legend
  \node[lblcap] at (0.3,5.2) {CONFIDENCE};
  \foreach \c/\t/\y in {ggmRecord/official{,} recent/4.88, ggmSource/older report/4.6, ggmAnalyst/similar country/4.32, red!65!black/guess/4.04} {
    \fill[\c] (0.42,\y) circle (0.07);
    \node[font=\sffamily\scriptsize, text=ggmMuted, anchor=west, inner sep=0] at (0.56,\y) {\t};
  }
  % --- suggestion card ---
  \filldraw[fill=white, draw=ggmAgent!50, line width=0.6pt, rounded corners=1.5mm, ggm soft shadow] (3.8,4.05) rectangle (7.45,10.3);
  \fill[ggmAgent, rounded corners=1.5mm] (3.8,9.85) rectangle (7.45,10.3);
  \fill[ggmAgent] (3.8,9.85) rectangle (7.45,10.05);
  \pic[scale=0.3] at (4.04,10.075) {agent badge};
  \node[font=\sffamily\scriptsize\bfseries, text=white, anchor=west] at (4.2,10.075) {Data agent suggests};
  \node[font=\sffamily\footnotesize\bfseries, text=ggmInk, anchor=west, inner sep=0.5mm] at (3.95,9.5) {Reference price};
  \node[font=\sffamily\scriptsize, text=ggmMuted, anchor=west, inner sep=0.5mm] at (3.95,9.17) {Gulf node, base year};
  \node[lblcap] at (3.95,8.8) {PROPOSED VALUE};
  \draw[line width=2pt, draw=ggmAgent!70, line cap=round] (4.05,8.55) -- (5.35,8.55);
  \node[lblcap] at (3.95,8.2) {WHY};
  \node[font=\sffamily\scriptsize, text=ggmInk, anchor=north west, text width=3.35cm, inner sep=0.5mm, align=left] at (3.95,8.05)
    {The official domestic price, the latest published.};
  \node[lblcap] at (3.95,7.12) {SOURCE};
  \node[lblbody] at (3.95,6.85) {National regulator};
  \node[lblcap] at (3.95,6.48) {CONFIDENCE};
  \node[ggm pill=ggmRecord, font=\sffamily\scriptsize\bfseries, inner ysep=0.7mm, anchor=west] at (3.98,6.16) {official, recent};
  \node[lblcap] at (3.95,5.76) {ALTERNATIVES};
  \draw[ph] (4.07,5.5) -- (5.5,5.5); \draw[ph] (4.07,5.25) -- (5.1,5.25);
  \tikzset{btn/.style={rounded corners=0.8mm, font=\sffamily\scriptsize\bfseries, minimum width=1.05cm, inner ysep=0.9mm}}
  \node[btn, draw=ggmRecord, text=ggmRecord] at (4.45,4.5) {Accept};
  \node[btn, draw=ggmMuted, text=ggmModel] at (5.63,4.5) {Edit};
  \node[btn, fill=ggmAnalyst, text=white] (reply) at (6.81,4.5) {Reply};
  % ---------- the analyst and the reply ----------
  \pic[scale=1.5] at (0.72,1.35) {person};
  \node[draw=ggmAnalyst, fill=ggmAnalyst!8, line width=0.7pt, rounded corners=2mm, text width=5.35cm, align=left,
        font=\sffamily\footnotesize, text=ggmInk, inner sep=2mm, anchor=west] (bubble) at (1.85,1.4)
    {\textbf{\color{ggmAnalyst!85!black}Reply:} This price is subsidised. It is far below what consumers would pay, so it is not the willingness to pay the model needs.};
  \fill[ggmAnalyst!8] (1.87,1.22) -- (1.55,1.33) -- (1.87,1.5) -- cycle;
  \draw[ggmAnalyst, line width=0.7pt] (1.87,1.22) -- (1.55,1.33) -- (1.87,1.5);
  \draw[ggm flow=ggmAnalyst, dashed, line width=0.8pt] (reply.south |- bubble.north) -- (reply.south);
  % ---------- where the answer goes ----------
  \coordinate (hub) at (8.05,1.4);
  \draw[ggmAnalyst, line width=1.1pt] (bubble.east) -- (hub);
  \draw[ggmAnalyst, line width=1.1pt] (hub) -- (8.05,9.6);
  \fill[ggmAnalyst] (hub) circle (0.07);
  \draw[ggm flow=ggmAnalyst] (8.05,9.6) -- (8.33,9.6);
  \draw[ggm flow=ggmAnalyst] (8.05,7.6) -- (8.33,7.6);
  \draw[ggm flow=ggmAnalyst] (8.05,4.0) -- (8.33,4.0);
  \node[font=\sffamily\scriptsize\bfseries, text=ggmAnalyst!85!black, rotate=90, fill=white, inner sep=0.6mm] at (8.05,5.75) {the reply goes to};
  % (a) back to the agent
  \filldraw[fill=white, draw=ggmAgent!45, rounded corners=2mm, ggm soft shadow] (8.35,9.0) rectangle (12.1,10.9);
  \pic[scale=0.55] at (8.8,10.45) {agent badge};
  \node[font=\sffamily\footnotesize\bfseries, text=ggmAgent, anchor=west] at (9.12,10.45) {Back to the agent};
  \node[font=\sffamily\scriptsize, text=ggmInk, anchor=north west, text width=3.35cm, inner sep=0, align=left] at (8.55,10.02)
    {The data agent revises the value and suggests again.};
  % (b) provenance ledger
  \filldraw[fill=white, draw=ggmRecord!45, rounded corners=2mm, ggm soft shadow] (8.35,5.95) rectangle (12.1,8.55);
  \pic[scale=0.6] at (8.8,8.07) {database};
  \node[font=\sffamily\footnotesize\bfseries, text=ggmRecord!85!black, anchor=west] at (9.12,8.07) {Provenance ledger};
  \begin{scope}[font=\sffamily\scriptsize, every node/.style={anchor=base west, inner sep=0}]
    \fill[ggmRecord!12] (8.5,7.15) rectangle (11.95,7.5);
    \foreach \x/\t in {8.6/Value, 9.55/Decision, 10.85/Why} \node[text=ggmRecord!80!black] at (\x,7.25) {\bfseries\t};
    \node at (8.6,6.87) {price}; \node[text=ggmAnalyst!85!black] at (9.55,6.87) {replied}; \node at (10.85,6.87) {subsidy};
    \draw[ggmRecord!25] (8.5,6.77) -- (11.95,6.77);
    \node at (8.6,6.51) {price}; \node[text=ggmRecord!85!black] at (9.55,6.51) {accepted}; \node at (10.85,6.51) {--};
    \node[text=ggmMuted] at (8.6,6.13) {who decided, and why};
  \end{scope}
  % (c) lesson
  \filldraw[fill=ggmRecord!7, draw=ggmRecord, line width=0.7pt, rounded corners=2mm, ggm soft shadow] (8.35,0.2) rectangle (12.1,5.5);
  \pic[scale=0.6] at (8.8,5.03) {doc=ggmRecord};
  \node[font=\sffamily\footnotesize\bfseries, text=ggmRecord!85!black, anchor=west] at (9.12,5.05) {New lesson};
  \node[font=\sffamily\scriptsize\itshape, text=ggmInk, anchor=north west, text width=3.35cm, inner sep=0, align=left] at (8.55,4.6)
    {Official prices where domestic gas is subsidised are not the willingness to pay. Flag them; propose alternatives.};
  \node[font=\sffamily\scriptsize, text=ggmMuted, anchor=west, inner sep=0] at (8.55,2.55) {Scope: all reference prices};
  \node[font=\sffamily\scriptsize, text=ggmMuted, anchor=west, inner sep=0] at (8.55,2.2) {Drawn from: the reply};
  \pic[scale=0.8] at (8.72,1.62) {check};
  \node[font=\sffamily\scriptsize\bfseries, text=ggmAnalyst!85!black, anchor=west, inner sep=0] at (8.97,1.62) {Analyst approved};
  \node[font=\sffamily\scriptsize, text=ggmMuted, anchor=north west, text width=3.35cm, inner sep=0, align=left] at (8.55,1.2)
    {Kept in a plain document the analysts can edit.};
  % revised suggestion goes back to the screen
  \draw[ggm flow=ggmAgent, line width=1pt, rounded corners=2mm] (8.35,10.45) -- (7.95,10.45) -- (7.95,10.075) -- (7.47,10.075);
  % lesson retrieved later
  \draw[ggm flow=ggmRecord, dashed, line width=1pt, rounded corners=3mm] (12.1,3.2) -- (12.4,3.2) -- (12.4,9.25) -- (12.12,9.25);
  \node[font=\sffamily\scriptsize\bfseries, text=ggmRecord!85!black, rotate=-90, fill=white, inner sep=0.6mm] at (12.4,6.5) {retrieved for later suggestions};
\end{tikzpicture}

\caption{A suggestion on the analyst's screen and what happens to the answer. The agent proposes a change with its reasons, source and confidence. The analyst accepts, edits or replies; the reply goes back to the agent, the decision goes to the provenance ledger, and a general rule drawn from it is kept as a lesson for later work. The screen is a sketch of a possible interface, and the example follows Section 5.8.}
\label{fig:review}
\end{figure}

\textbf{What a suggestion contains.} Next to the proposed values, the agent states why it proposes them, where they come from and how confident it is.
Confidence is given in the classes the developers proposed: taken from a recent official source, from an older report, inferred from a similar country and which similarity, or a guess.
Where there are reasonable alternatives, such as two outlooks that differ, the agent shows them side by side.
The classes let the analyst review selectively, starting with the guesses, instead of checking every value with equal care.

\textbf{Three answers.} The analyst can accept the suggestion, edit the values, or reply in plain language, for example that a price is too low because the country subsidises domestic gas.
A reply sends the suggestion back to the agent, which revises it with the new information and suggests again.
The developers proposed a similar exchange with whole workbooks: the analyst returns an edited workbook with a column of overrides and notes, and the agent takes it into the next round.
Either way, every decision is written to the provenance ledger, so that each value records who proposed it, who accepted it and on what grounds.

\textbf{Learning from corrections.} A reply often contains a rule that holds beyond the case at hand.
That official prices in countries with subsidised domestic consumption are far below the willingness to pay the model needs is one the developers named, and one that LLMs handle badly.
When a reply contains such a rule, the agent proposes it as a lesson, worded as an instruction, and the analyst approves or rewrites it.
Approved lessons are kept in a plain document and retrieved into the agent's context whenever they are relevant, in the same way that retrieval supplies an LLM with documentation \citep{lewis2020retrieval}.
Agents that keep a written reflection on each failure and use it on the next attempt have been shown to improve without any retraining \citep{shinn2023reflexion}, and the lessons apply the same idea across a whole study, and from one study to the next.

This is how the system would improve with each working hour.
Nothing in the language model changes; what grows is a document of rules that the analysts can read, edit and remove.
Much of the expertise in GGM's data preparation already has this form: rules such as using marketed rather than gross production, or how a country's total is split across its nodes, which the documentation records for a few countries \citep[pp.~24--25]{egging2019global} and the developers otherwise carry in their heads and in earlier spreadsheets.
Each correction thus also records a piece of that expertise.
And each correction becomes a test case: a suggestion the analyst rejected, with the reason, is something a later version of the system should get right.

The lessons carry their own risk.
A wrong lesson would be applied silently to every later case it seems to fit, so lessons need the analyst's approval, a stated scope, and a record of where they have been applied.
A second risk is that an analyst who accepts most suggestions starts to accept them without reading; Chapter 6 returns to this.
